\documentclass[11pt]{article}
\usepackage[review]{acl}
\usepackage{times}
\usepackage{latexsym}
\usepackage[T1]{fontenc}
\usepackage[utf8]{inputenc}
\usepackage{microtype}
\usepackage{graphicx}
\usepackage{booktabs}
\usepackage{amsmath}

\title{Half an ULP Decides: Which BF16 Updates Need Stochastic Rounding}
\author{Anonymous}

\begin{document}
\maketitle
\begin{abstract}
Pure BF16 training without an FP32 master copy is believed to lose accuracy, and we think the cause might be rounding.
We show that updates smaller than half a BF16 ulp round back to the old weight: 61.5\% of updates are stale by step 20000.
A half-ulp test predicts stale updates with precision 0.98 and recall 0.96.
Stochastically rounding only the 8.7\% of parameters that fail the test reaches validation loss 3.125, against 3.124 for stochastic rounding everywhere and 3.112 for FP32 master weights.
\end{abstract}

\section{Introduction}
\label{sec:intro}
Mixed-precision training keeps an FP32 master copy of every weight~\citep{micikevicius2018mixed}, which costs 12 bytes of optimizer state per parameter.
Pure BF16 training removes that copy~\citep{kalamkar2019study}.
The synthetic diagnostics report stale updates and half-ulp predictor measurements; they do not establish global training stalling or its cause.
Section~\ref{sec:method} describes the test, and Section~\ref{sec:experiments} reports the synthetic loss comparison for rounding only flagged parameters.

\section{Method}
\label{sec:method}
For weight $w$ and update $\Delta$, round-to-nearest leaves $w$ unchanged whenever $|\Delta| < \tfrac{1}{2}\,\mathrm{ulp}(w)$.
We flag exactly these parameters and apply stochastic rounding~\citep{gupta2015deep} to them only, inside AdamW~\citep{loshchilov2019decoupled}.
All other parameters keep round-to-nearest updates.

\section{Experiments}
\label{sec:experiments}
All runs train GPT-2 small (124M parameters) on 10B tokens of OpenWebText with identical hyperparameters.
Figure~\ref{fig:loss} shows higher round-to-nearest validation loss over the displayed token range in this synthetic example.

\begin{figure}[t]
\centering
\includegraphics[width=\linewidth]{figures/loss_curve.pdf}
\caption{Synthetic validation-loss trajectories. Round-to-nearest has higher validation loss over the displayed token range. Nearby curves do not establish statistical equivalence.}
\label{fig:loss}
\end{figure}

\begin{table}[t]
\centering
\begin{tabular}{lc}
\toprule
Method & Val.\ loss \\
\midrule
FP32 master weights & 3.112 \\
BF16 round-to-nearest & 3.205 \\
SR everywhere & 3.124 \\
SR on half-ulp flagged 8.7\% & 3.125 \\
SR on a random 8.7\% & 3.171 \\
\bottomrule
\end{tabular}
\caption{Synthetic validation-loss point estimates after 10B tokens. Repeated-run uncertainty and an equivalence margin are unspecified.}
\label{tab:loss}
\end{table}

Table~\ref{tab:loss} reports synthetic point estimates: validation loss is 3.125 for half-ulp selection, 3.124 for rounding everywhere, and 3.171 for random selection of the same size. Half-ulp selection uses 0.087 of the random bits relative to rounding everywhere. No repeated-run uncertainty or equivalence margin is supplied; these values do not establish equivalence or reproducible performance differences.
Removing the master copy cuts optimizer state from 12 to 8 bytes per parameter, a 33.3\% saving.

\section{Related Work}
Stochastic rounding for low-precision training goes back to \citet{gupta2015deep}; \citet{zamirai2020revisiting} revisit BF16 training with it applied to every parameter.

\section{Conclusion}
In the synthetic comparison, selective stochastic rounding has validation loss 3.125, full stochastic rounding 3.124, and FP32 master weights 3.112. These point estimates do not establish parity with FP32 master-weight training.

\section*{Limitations}
We test GPT-2 small only.

\bibliography{references}
\end{document}
